\SourceLeaf{140}
\begin{center}
NACHLASS UND BRIEFWECHSEL.
\end{center}

\begin{center}
[I.]
\end{center}

\begin{center}
[AUFGABE AUS DER WAHRSCHEINLICHKEITSRECHNUNG.]
\end{center}

\bigskip
\hrule
\bigskip

Es sei $\mu$ die Wahrscheinlichkeit eines bestimmten Erfolges $E$ aus einem
einfachen Versuche, mithin $1-\mu$ die Wahrscheinlichkeit des Ausbleibens dieses
Erfolges.

Die Wahrscheinlichkeit, dass unter $n$ von einander unabhängigen Versuchen $k$
den Erfolg $E$ geben, ist

\[
\frac{n(n-1)(n-2)....(n-k+1)}{1.2.3....k}(1-\mu)^{n-k}\mu^k = \varphi k.
\]

\begin{center}
\small\setlength{\tabcolsep}{3pt}
\begin{tabular}{lllllcl}
Der & mittlere & Werth & von & $k$ & ist & $n\mu$, \\
» & » & » & » & $k(k-1)$ & » & $n(n-1)\mu\mu$, \\
» & » & » & » & $k(k-1)(k-2)$ & » & $n(n-1)(n-2)\mu^3$, \\
\multicolumn{7}{c}{u.\,s.\,w.}\\
» & » & » & » & $(k-n\mu)^2$ & » & $n\mu(1-\mu)$. \\
\end{tabular}
\end{center}

\bigskip
\hrule
